% BEGIN OPTIMIZED PACKED PI-LOGSTAR EVALUATION
\subsection{Concrete Optimization of \ourprotocol}\label{sec:packed-eval}
\input{plots/packed/packed-values.tex}

We now evaluate an optimized implementation of the main construction.
This subsection supplements the preceding development measurements.
It uses a new executable and fresh experiments at every power of two from
$n=2^{10}$ through $n=2^{20}$, where $n$ is the number of keys in each list.
Both protocols return the same stable, XOR-shared gather permutation of $2n$
inputs. The baseline is a Batcher bitonic merge network on the same GMW backend,
with balanced comparisons and projection of the final layer onto required output bits.
The optimized implementation reduces online communication by
\packedSmallReduction\% at $2^{10}$ and \packedLargeReduction\% at $2^{20}$
relative to Batcher. We report preprocessing separately throughout.

\paragraph{Platform and measurement procedure.}
The platform is the Lenovo 83DF laptop described above, with an Intel Core
i9-14900HX and approximately $32$\,GiB of physical memory.
Ubuntu under WSL2 exposes $32$ logical CPUs and approximately $15$\,GiB of RAM.
We compile C++20 in Release mode with GCC 13.3 and CMake 3.28.3,
using \texttt{-O3} and \texttt{-march=native}.
Local experiments execute both parties as coroutines on one OS thread.
Two concurrent correlation batches overlap coroutine work; they do not create
two cryptographic worker threads. Each TCP party runs in a separate process
with two inherited Boost.Asio I/O workers. CPU affinity and frequency are not fixed.

Every local size has three fresh executions per protocol. We alternate which
protocol runs first across repetitions. Both protocols use $32$-bit keys,
$2^{20}$-entry correlation batches, and the same synthetic input seed within
each paired trial. Private cryptographic randomness comes from independent
OS-seeded PRNGs. All preprocessing uses the real inherited OT/OLE and permutation
protocols; no dealer, mock, or idealized substitute is timed.
We report medians and show observed minimum--maximum timing ranges.
Setup, output opening, and correctness checks are excluded from phase timings.
Each measured output is checked against an independent stable merge.
Communication sums both parties' sent bytes once, includes coproto framing,
and excludes TCP/IP headers.\footnote{Local Batcher execution charges a
24-byte session registration per party to the online phase. TCP initializes
that session in the excluded configuration handshake. We retain the measured
transport counters without normalization.} MiB denotes $2^{20}$ bytes.
We record peak process RSS and sample process swap every $0.5$ seconds.

\paragraph{Implementation and public parameter selection.}
The five protocol components and shared permutation header comprise
\packedLibraryLoc\ physical lines of
C++ source and headers, including comments and blank lines.
The benchmark and focused C++ tests comprise another \packedTestLoc\ lines.
These counts exclude inherited cryptographic dependencies and other repository
components. The optimized path applies to power-of-two inputs with one partition
and terminal block size at most $16$. General padded and deeper recursive schedules
remain implemented separately.

For each input size, we independently evaluate public circuit costs for
$m\in\{2,4,8,16\}$. The objective is the padded GMW communication plus the
block-dependent permutation payload. Terms common to all candidates are omitted
from this selection objective. The minimizing block size is $m=4$ at every
measured size; each base merge therefore combines two four-row blocks.
These are public parameter choices, independent of the keys and timing outcomes.
The artifact records all candidate costs and the selected schedule.
This fixed one-partition specialization targets concrete costs; we make no
new asymptotic claim for a fixed block size as $n$ grows.

\paragraph{Compressed blocks and parallel base merges.}
A block carries its keys and one original block identifier.
Each row's original index is recovered from that identifier and its public
offset, avoiding repeated indices during block permutation and broadcast.
The representative merge also avoids a duplicate 32-bit payload index.
The broadcast reduces groups of eight blocks, scans group endpoints with a
Sklansky network, and distributes the resulting prefixes within groups.
For $L>8$ blocks, its depth is $\log_2 L+3$, instead of $2\log_2 L-1$.
This reduces depth at a modest increase in broadcast work.

Each terminal merge compares all cross-pairs in parallel.
The comparisons use the key and original source bit to resolve cross-list ties.
Sortedness makes the comparison results monotone; adjacent XORs produce one-hot
insertion positions. Public offsets and source choices are then accumulated
with XORs, and validity bits require one product per candidate.
This replaces repeated comparator/swap layers with one comparison batch.
Copied rows with no predecessor from the original block are inactive, so the
base merge also implements the lower interval mask.
Upper masks compare keys and source bits, using the sorted block order to
resolve equal same-source boundaries. We preserve each real row exactly once.

\paragraph{Direct ranks and stable extraction.}
An active row's global rank is the sum of the two within-list block-start
indices and its position in the local merged block.
The interval masks ensure that a later opposite-source block cannot contribute
an earlier key. When the opposite predecessor is absent, its dummy start and
cross-comparison contributions are zero. Two additions per block compute rank
prefixes for the two halves of the output, avoiding per-row arithmetic conversion
and extraction OTs.

We correct the inherited permutation sampler before using it here.
The previous sampler reduced 32-bit random words modulo the remaining length,
introducing bias. We use \texttt{std::shuffle} with the cryptographic PRNG's
64-bit generator interface and unbiased bounded sampling.
The correction also applies to block permutations.
The preceding development measurements predate this correction; all experiments
in this subsection use the corrected sampler.

We jointly shuffle the shared validity flags, ranks, and output indices using
fresh permutation correlations. We then open the shuffled flags and only the
ranks at valid positions, and place the still-shared indices at those ranks.
For either semi-honest party, the other party's private random permutation hides
the original positions. In the hybrid with ideal primitives and uniform randomness,
the opened transcript is a uniform placement of the
distinct labels $0,\ldots,2n-1$ among $4n$ positions; its distribution depends
only on public dimensions. Inactive ranks and output indices remain shared.
The artifact gives the corresponding simulation argument under the inherited
primitive assumptions. Correlations are consumed once.

These optimizations pass $8{,}642$ independent plaintext cases with direct-rank
checks, $72$ additional real-cryptography cases with random XOR shares, and
separate extraction, prefix, and permutation tests. Validation includes duplicate and maximum
keys, wide keys, padding, deeper recursion, and rejection of correlation reuse.
An exhaustive four-position example also checks the input independence of the
shuffled-rank transcript; it supplements the argument rather than replacing it.

\paragraph{Communication and local execution.}
\Cref{fig:packed-communication} shows the complete sweep.
The optimized protocol uses fewer online bytes already at the smallest sampled
size, $2^{10}$: \packedSmallPi\ versus \packedSmallBatcher\,MiB.
At $2^{20}$, it uses \packedLargePi\ versus \packedLargeBatcher\,MiB,
a \packedLargeReduction\% reduction. Its online traffic is also
\packedReductionOld\% below the previous implementation's $2{,}125.159$\,MiB.
The first sampled size already favors the optimized protocol; we do not infer
an exact break-even size below the evaluated range.

\begin{figure}[!htbp]
\centering
\includegraphics[width=\linewidth]{plots/packed/packed-communication.pdf}
\caption{Verified online communication at every power of two from $2^{10}$
through $2^{20}$ keys per list. The right panel divides optimized
\ourprotocol\ bytes by Batcher bytes; the dashed line denotes equality.
Connecting lines guide the eye.}
\label{fig:packed-communication}
\end{figure}

\input{plots/packed/packed-table.tex}

At $2^{20}$, the local medians for \ourprotocol\ and Batcher are
\packedLargeOnlinePi\ and \packedLargeOnlineBatcher\,s online, and
\packedLargeOfflinePi\ and \packedLargeOfflineBatcher\,s preprocessing.
Peak RSS is \packedLargeRssPi\ and \packedLargeRssBatcher\,GiB, respectively.
\packedPagingStatement{}
The online ranges overlap, so these medians do not establish a consistent
local speed advantage.

\begin{figure}[!htbp]
\centering
\includegraphics[width=\linewidth]{plots/packed/packed-time-depth.pdf}
\caption{Local online time and preprocessing time, with both parties on one
OS thread. Markers are three-run medians; error bars show observed ranges.}
\label{fig:packed-runtime}
\end{figure}

\paragraph{Preprocessing and depth.}
The byte savings above concern the online phase.
At $2^{20}$, \ourprotocol\ preprocessing sends \packedLargeOfflineBytesPi\,MiB.
Batcher sends \packedLargeOfflineBytesBatcher\,MiB.
The corresponding online depth bounds are \packedLargeDepthPi\ and
\packedLargeDepthBatcher. We count GMW AND layers and nine additional one-way
steps for block permutation and shuffled extraction; preprocessing is excluded.
These are dependency-depth bounds, not packet counts.
The optimized construction uses more preprocessing bandwidth and greater online
depth, even where it reduces online communication and execution time.

\paragraph{Network sensitivity.}
We use isolated Linux network namespaces connected by a virtual Ethernet pair.
Each direction receives a \texttt{tc netem} rate limit and half the configured
RTT. The profiles are LAN ($1$\,Gbit/s, $0.2$\,ms), WAN ($100$\,Mbit/s,
$40$\,ms), and slow WAN ($10$\,Mbit/s, $40$\,ms).
These are emulated links on one laptop. We take three paired trials per profile
at $2^{12}$ and $2^{16}$, and one paired WAN trial at $2^{20}$.
TCP phase time is the slower party's time.
\input{plots/packed/packed-network-text.tex}

\begin{figure}[!htbp]
\centering
\includegraphics[width=\linewidth]{plots/packed/packed-network.pdf}
\caption{Measured TCP online time and offline-plus-online time.
Bars at $2^{12}$ and $2^{16}$ are three-run medians with observed ranges;
the $2^{20}$ WAN comparison is one execution per protocol.
Rates are full-duplex limits per direction.}
\label{fig:packed-network}
\end{figure}

\paragraph{Batch sensitivity and reproducibility.}
A separate matched pilot compares correlation batches of $2^{20}$ and $2^{22}$
at $2^{14}$ and $2^{16}$. \packedBatchStatement{}
We retain $2^{20}$ for the main sweep
and all network comparisons, identically for both protocols.
Online byte counts do not depend on this preprocessing batch choice.
The artifact contains source, all candidate schedules, raw paired JSONL records,
memory measurements, plotting scripts, and the executable hash.
\clearpage
% END OPTIMIZED PACKED PI-LOGSTAR EVALUATION
